\providecommand{\kBKM}{\kappa_{\mathrm{BKM}}}
\providecommand{\kch}{\kappa_{\mathrm{ch}}}
\providecommand{\kcat}{\kappa_{\mathrm{cat}}}
\providecommand{\kfib}{\kappa_{\mathrm{fiber}}}

\section*{The Nonexistent Bryan, Oberdieck, and Pandharipande Theorem}

\subsection*{Bibliographic Fact}

The reference string
\[
\text{``Bryan, Oberdieck, and Pandharipande 2019, Theorem 1''}
\]
with identifier \texttt{arXiv:1903.07729} belongs to the
materials-science paper
\emph{Constitutive modelling of mechanically induced martensitic
transformations: Prediction of transformation surfaces} by
D.~de~Bortoli, F.~Adziman, E.~A.~de~Souza~Neto, and
F.~M.~Andrade~Pires.  Its subject is constitutive mechanics of
martensitic transformations, outside reduced Gromov-Witten theory,
reduced Donaldson-Thomas theory, CHL quotients, Igusa forms, and
Borcherds products.

The relevant surrounding literature has the following precise scope.

\begin{itemize}
\item Oberdieck and Pandharipande, \texttt{arXiv:1411.1514},
conjecture the $K3\times E$ reduced Gromov-Witten formula in terms of
the Igusa cusp form $\chi_{10}$ and give primitive descendent evidence.

\item Bryan, \texttt{arXiv:1504.02920}, computes reduced
Donaldson-Thomas partition functions for primitive classes of square
$-2$ and $0$ on $K3\times E$ via the topological vertex.

\item Oberdieck and Pixton, \texttt{arXiv:1706.10100},
\emph{Inventiones Mathematicae} \textbf{222} (2020), 667 to 759, solve
the reduced Gromov-Witten theory of $S\times E$ for all curve classes
primitive in the K3 factor and deduce the Igusa cusp-form conjecture.
This is the untwisted $N=1$ Gromov-Witten theorem.

\item Bryan and Oberdieck, \texttt{arXiv:1811.06102}, formulate the
CHL Donaldson-Thomas conjecture for elliptic CHL models
$[K3\times E/\mathbb{Z}_N]$ at $N\in\{2,3,4,6\}$.  Their proved
results are topological-vertex base cases at $N=2$, in primitive
classes of square $-2$ and $0$.

\item Bryan, Oberdieck, Pandharipande, and Yin,
\texttt{arXiv:1506.00841}, study curve counting on abelian surfaces
and abelian threefolds.  That four-author paper belongs to a different
geometric class.
\end{itemize}

Thus the literature chain contains an $N=1$ Gromov-Witten theorem,
primitive reduced Donaldson-Thomas checks, and a CHL
Donaldson-Thomas conjecture with $N=2$ base cases.  The all-order,
all-$N$ three-author 2019 theorem is absent from this chain:
\[
 Z^{\mathrm{red}}_{\mathrm{GW}}
 =
 Z^{\mathrm{red}}_{\mathrm{DT}}
 =
 \Phi_N^{-1}
\]
for all $N\in\{1,2,3,4,6\}$ on crepant resolutions of
$[K3\times E/\mathbb{Z}_N]$.

\subsection*{Primitive CHL/BKM Constants}

Let
\[
 \mathcal{N}_{\mathrm{CHL}}=\{1,2,3,4,6\}.
\]
For $N\in\mathcal{N}_{\mathrm{CHL}}$, let $\Phi_N$ denote the primitive
CHL/BKM denominator in the Vol~III normalisation, and let $c_N(0)$ be
the $\zeta^0q^0$ coefficient of the weak Jacobi form entering the
Borcherds lift.  The Borcherds weight theorem gives
\[
 \kBKM(\Phi_N)=\frac{c_N(0)}{2}.
\]
The primitive CHL row is
\[
\begin{array}{c|ccccc}
N & 1 & 2 & 3 & 4 & 6\\
\hline
c_N(0) & 10 & 8 & 6 & 4 & 2\\
\kBKM(\Phi_N) & 5 & 4 & 3 & 2 & 1
\end{array}
\]
with the $N=1$ entry
\[
 \kBKM(\mathfrak g_{\Delta_5})
 =
 \operatorname{wt}(\Delta_5)
 =
 \frac{c_1(0)}{2}
 =
 5.
\]
This table is an automorphic statement: Borcherds' product formula
reads the weight from the constant Fourier coefficient of the input
Jacobi form.  The reduced Donaldson-Thomas interpretation of CHL
quotients is a separate geometric problem.

\subsection*{Gritsenko-Clery Catalogue}

The Gritsenko-Clery diagonal-divisor catalogue is a separate eight-row
paramodular atlas.  Its rows are indexed by triples
\[
 (t,N_{\mathrm{lvl}};k)\in
 \{(1,1;5),(2,1;2),(1,2;3),(3,1;1),
 (1,3;2),(4,1;1/2),(1,4;3/2),(2,2;1)\}.
\]
The row weights are
\[
 (5,2,3,1,2,1/2,3/2,1),
\]
and the corresponding twice-weight tuple is
\[
 (10,4,6,2,4,1,3,2).
\]
The primitive CHL tuple is $(10,8,6,4,2)$, a different tuple.  The two
atlases share the single entry $(1,1;5)$, the $\Delta_5$ row.  Any
comparison between a CHL denominator at $N\in\{2,3,4,6\}$ and a
Gritsenko-Clery row requires a stated row map, multiplier, level, input
Jacobi form, and denominator normalisation.

\subsection*{\texorpdfstring{$K3\times E$}{K3 x E} Invariants}

For the compact product $X=K3\times E$, the Hodge supertrace gives
\[
 \kch(X)
 =
 \sum_{q=0}^3(-1)^qh^{0,q}(X)
 =
 1-1+1-1
 =
 0.
\]
The categorical Euler characteristic is multiplicative:
\[
 \kcat(X)
 =
 \chi(\mathcal{O}_{K3})\,\chi(\mathcal{O}_E)
 =
 2\cdot 0
 =
 0.
\]
The rank-additive Heisenberg specialisation has
\[
 \kch^{\mathrm{Heis}}=3,
\]
and the K3 fibre contributes the Mukai-lattice rank
\[
 \kfib(K3)
 =
 \operatorname{rank}\widetilde H(K3,\mathbb{Z})
 =
 24.
\]
The Borcherds denominator contributes the independent automorphic
constant
\[
 \kBKM(\mathfrak g_{\Delta_5})=5.
\]
The four values
\[
 \kch(K3\times E)=0,\qquad
 \kcat(K3\times E)=0,\qquad
 \kBKM(\mathfrak g_{\Delta_5})=5,\qquad
 \kfib(K3)=24
\]
come from distinct constructions.  The additional Heisenberg reading
$\kch^{\mathrm{Heis}}=3$ belongs to the rank-additive
Mukai-Heisenberg specialisation.  The compact Hodge supertrace is the
separate value $\kch(K3\times E)=0$.
The numerical coincidence
\[
 5=3+2
\]
mixes the Heisenberg value with the fibre value
$\chi(\mathcal{O}_{K3})=2$ and therefore carries zero functorial force on
$K3\times E$.

\subsection*{Mathematical Consequence}

The automorphic identity
\[
 \kBKM(\Phi_N)=c_N(0)/2
\]
is theorem-level on the Borcherds side for the primitive CHL constants
displayed above.  The geometric equality of reduced
Gromov-Witten and reduced Donaldson-Thomas partition functions on all
CHL quotients at $N\in\{2,3,4,6\}$ remains a separate proof obligation:
one must supply the reduced CHL Gromov-Witten theory, the reduced CHL
Donaldson-Thomas theory in all primitive classes, the orbifold or
crepant-resolution comparison, and the reduced GW/DT correspondence in
the same normalisation.

\subsection*{References}

\begin{itemize}
\item R.~E.~Borcherds, \emph{Automorphic forms with singularities on
Grassmannians}, \emph{Inventiones Mathematicae} \textbf{132} (1998),
491 to 562.
\item V.~Gritsenko, \emph{Arithmetical lifting and its applications}.
\item V.~Gritsenko and V.~Nikulin, Siegel automorphic form corrections
of Lorentzian Kac-Moody algebras.
\item F.~Clery and V.~Gritsenko, diagonal-divisor paramodular forms.
\item G.~Oberdieck and R.~Pandharipande, \emph{Curve counting on
$K3\times E$, the Igusa cusp form $\chi_{10}$, and descendent
integration}, \texttt{arXiv:1411.1514}.
\item J.~Bryan, \emph{The Donaldson-Thomas theory of $K3\times E$ via
the topological vertex}, \texttt{arXiv:1504.02920}.
\item G.~Oberdieck and A.~Pixton, \emph{Holomorphic anomaly equations
and the Igusa cusp form conjecture}, \texttt{arXiv:1706.10100}.
\item J.~Bryan and G.~Oberdieck, \emph{CHL Calabi-Yau threefolds:
curve counting, Mathieu moonshine and Siegel modular forms},
\texttt{arXiv:1811.06102}.
\item J.~Bryan, G.~Oberdieck, R.~Pandharipande, and Q.~Yin,
\emph{Curve counting on abelian surfaces and threefolds},
\texttt{arXiv:1506.00841}.
\end{itemize}
